\section{Design decisions and trade-offs}\label{sec:decisions}

Each decision below is a short story: one agent found a problem, and another acted on it.

\subsection{PSRAM: from WROOM-32E to WROVER-E}\label{sec:psram}

The firmware research found that Espressif's audio framework says Bluetooth and Wi-Fi ``can not coexist without PSRAM''. With both stacks running, about \qtyrange{100}{150}{\kilo\byte} of internal heap remains, while one full LVGL frame alone needs \qty{115}{\kilo\byte} \src{research/firmware_and_tools.md}. The lead switched to the ESP32-WROVER-E-N8R8 while the electronics agent was still running. At quantity 10 it was even cheaper than the WROOM (\$3.73 against \$4.04). It cost \qty{6}{\milli\metre} of module length and GPIO16/17, which left \emph{no spare GPIO} on the board \src{research/electronics.md}.

\subsection{\texorpdfstring{\qty{4}{\ohm}}{4 Ω} or \texorpdfstring{\qty{8}{\ohm}}{8 Ω} woofers, and the amplifier current limit}\label{sec:impedance}

The TAS5805M trips at \qty{5}{\ampere}. A \qty{4}{\ohm} woofer on a \qtyrange{19}{20}{\volt} rail reaches \qtyrange{4.75}{5}{\ampere} at clipping. The driver agent re-ran its alignment with the amplifier modelled as $V_\text{peak}\le\min(\qty{18}{\volt}, \qty{4.5}{\ampere}\cdot|Z|)$. Going to \qty{8}{\ohm} cost 1--\qty{3}{\decibel}: the best \qty{4}{\ohm} set was estimated at \qty{94.4}{\decibel} pink, against \qty{91.4}{\decibel} for the \qty{8}{\ohm} set B8 \src{research/drivers.md}. After the flip-flop of \cref{sec:process}, the lead chose \textbf{B8 with PVDD = \qty{18}{\volt}}. A clipped woofer then draws only \qty{2.28}{\ampere}.

\textbf{The problem then moved to the tweeter.} The \qty{4}{\ohm} OT19NC00-04 ($R_e=\qty{3.17}{\ohm}$) would draw \qty{5.2}{\ampere} at hard clipping, including ripple. The signal spec therefore made a tweeter limiter at $\le\qty{8}{\volt}$ peak mandatory \src{docs/spec_signal.md}. SPICE confirmed it independently: \qty{5.41}{\ampere} with a clipped tone, and $\le\qty{2.7}{\ampere}$ with the limiters (\cref{sec:spice}).

\subsection{The battery hunt and the Ø\texorpdfstring{\qty{95}{\milli\metre}}{95 mm} capsule}\label{sec:battery}

The user rejected spot-welded cells, so the plan became a pre-built 2S pack with protection on the board. The mechanical agent found that the geometry decided the rest \src{docs/mech_envelope.md}:
\begin{itemize}
  \item The \qty{48}{\milli\metre}-deep woofers, rear-mounted, leave only a \qtyrange{19}{23}{\milli\metre} slab behind their motors.
  \item At Ø\qty{90}{\milli\metre} that slab shrinks to about \qty{16}{\milli\metre}, which would need a pack no thicker than \qty{13}{\milli\metre}. No stocked pack was that thin with enough capacity.
\end{itemize}
The choice was a Gens ace \qty{5000}{\milli\ampere\hour} 2S RC pack: \qtyproduct{109 x 42 x 19.5}{\milli\metre}, \qty{37}{\watt\hour}, about \$31. It fits a Ø\qty{95}{\milli\metre} capsule with \qty{1}{\milli\metre} to spare and a \qty{10}{\percent} swell allowance. At a \qty{2.9}{\watt} ``moderate'' average it runs about \qty{12.8}{\hour}, against the \qty{8}{\hour} target. The pack has no protection circuit, no NTC (one is bonded on) and no published continuous rating, so it must be bench-tested at \qty{6}{\ampere}. The board's protection was designed for exactly this kind of pack.

\subsection{Removing the boost load-disconnect FET}\label{sec:disconnect}

TI's rules for the TPS61178's optional disconnect FET require the capacitance after the FET to be at most $10\times$ the capacitance before it. Here it would be about $40\times$ (\qty{1.5}{\milli\farad} against \qty{36}{\micro\farad}), and start-up would risk tripping the hiccup protection. The power-spec agent removed the FET \src{docs/spec_power.md}. PVDD now rests at VSYS$-\qty{0.5}{\volt}$ when the boost is off, with the amplifiers in shutdown drawing about \qty{30}{\micro\ampere}.

The release review later showed what the spec had assumed:
\begin{itemize}
  \item \textbf{Short protection depends on firmware.} A PVDD short is cleared only by the charger's \qty{9.3}{\ampere} battery over-current trip, which works only while firmware keeps it armed; a watchdog expiry disarms it (S1).
  \item \textbf{Inrush at pack insertion is not bounded:} \qtyrange{2.5}{9.7}{\ampere} through the ship FET (S2).
\end{itemize}
The first became the top item of the firmware contract (\cref{sec:firmware-req}). The second added a DNP gate--source footprint, C721, with a corrected fitting rule (N1, \cref{tab:firmware-req}).

\subsection{The ADR/FAULT pin}\label{sec:adr-fault}

The brief asked for one shared open-drain fault line from both amplifiers. The signal-spec agent found that the TAS5805M's pin 3 is \texttt{ADR/FAULT}: it reads the I\textsuperscript{2}C address resistor at power-up, then becomes a push-pull FAULTZ output. Tying the two pins together would have given \emph{both amplifiers the same address}. The fix keeps separate address resistors and merges the two FAULTZ signals with two 2N7002s in series. The merged AMP\_FAULT on GPIO36 is active-high, and the FETs load the address straps by only \qty{1.5}{\percent} \src{docs/spec_signal.md}.

\subsection{CH224Q back-driving the I\texorpdfstring{\textsuperscript{2}}{2}C bus}\label{sec:ch224q}

The PD controller shares the I\textsuperscript{2}C bus but is powered only from USB. Unpowered, its internal pull-ups become pull-downs and would drag the \qty{3.3}{\volt} bus to about \qty{2.2}{\volt}, below the input-high threshold. On battery the speaker would have lost control of its amplifiers. Two FETs with their gates on VBUS/2 now isolate the PD controller when USB is absent. They also set a firmware rule: never request \qty{28}{\volt} EPR, because the gates would see \qty{14}{\volt} \src{docs/spec_power.md}. The release review later found that the FETs' output capacitance may slow the bus edges, so I\textsuperscript{2}C runs at \qtyrange{100}{200}{\kilo\hertz} until the rise time is measured (S5).

\subsection{Smaller catches}\label{sec:smaller-catches}

\begin{itemize}[itemsep=1pt]
  \item The HY2213 balancers are gate drivers, not bleeders: each cell needs its own FET and resistors.
  \item ESP32 erratum: GPIO36/39 glitch when ADC1 starts, so the encoder switch moved to GPIO4.
  \item No JLC stock of JST-SH, so the UI link became a 20-pin FFC, with its own signal-integrity study (\cref{sec:spice}).
  \item Battery negative is not ground (low-side protector), so BATT\_N must be its own copper island.
  \item The charger's power-on charge current is \qty{2}{\ampere}, not the \qty{1}{\ampere} copied from its sibling chip (S3).
  \item Charging temperature is now sensed on the pack's own NTC, not on the board (ECO-002).
  \item An external power button was added through FFC pin 20. Firmware cannot see it, so it only wakes and resets; ``off'' is a long press of the encoder (N2).
\end{itemize}

\subsection{Change orders}\label{sec:eco}

The SPICE/PDN agent wrote its findings as an engineering change order rather than a report (\cref{tab:eco}) \src{docs/eco_001.md}. ECO-002 did the same for the release review; it is summarised in \cref{app:register}. Both are applied in schematic v0.2.

\begin{table}[htbp]
\centering\small
\caption{ECO-001, from the circuit simulations. Items 1--4 are applied in the schematic; 5--6 are firmware requirements.}
\label{tab:eco}
\begin{tabularx}{\linewidth}{@{}c >{\raggedright\arraybackslash}p{0.36\linewidth} >{\raggedright\arraybackslash}X@{}}
\toprule
\# & Change & Evidence \\
\midrule
1 & Boost compensation \qty{33}{\kilo\ohm}/\qty{47}{\nano\farad} $\to$ \qty{100}{\kilo\ohm}/\qty{15}{\nano\farad} & PVDD impedance \qtyrange{0.46}{0.57}{\ohm} at \qty{150}{\hertz} against a \qty{0.14}{\ohm} target; load-step dip \qty{0.97}{\volt} $\to$ \qty{0.34}{\volt} \\
2 & Fit the trim network R907/D901 for a firmware soft start & Start-up windup overshoots to \qty{19.4}{\volt} (typ.) or the \qty{21}{\volt} OVP \\
3 & SPI/PDM source resistors \qty{33}{\ohm} $\to$ \qty{68}{\ohm} & \qtyrange{0.7}{1.1}{\volt} overshoot on the FFC \\
4 & ESP32 EN capacitor \qty{1}{\micro\farad} $\to$ \qty{2.2}{\micro\farad} & Auto-reset fails if the host's DTR/RTS gap exceeds \qty{2.5}{\milli\second} \\
5 & Tweeter EQ: \qty{+2}{\decibel} high shelf (firmware) & LC filter droop \qty{-2.2}{\decibel} at \qty{20}{\kilo\hertz} \\
6 & Stop charging above \qty{45}{\degreeCelsius} (firmware) & The hardware window runs to \qty{60.8}{\degreeCelsius} \\
\bottomrule
\end{tabularx}
\end{table}
